% COMS 6998, Fall 2026 -- Project Proposal
% Adapted from the official ISCA 2026 IEEE template.
% Overleaf: upload the ZIP, select main.tex, and use pdfLaTeX.
% Local: latexmk -pdf main.tex
%
% Replace the guidance and bracketed fields with your project content.
% Use the five sections below and retain the supplied formatting.
% Main content: two pages, including figures, tables, and the team plan.
% References are excluded from the two pages.

\documentclass[conference,10pt,letterpaper]{IEEEtran}
\usepackage[T1]{fontenc}
\usepackage{cite}
\usepackage{amsmath,amssymb}
\usepackage{graphicx}
\usepackage{booktabs}
\usepackage{tabularx}
\usepackage{flushend}
\usepackage[hyphens]{url}
\usepackage[hidelinks]{hyperref}

\newcommand{\field}[1]{\textit{[#1]}}
\newcolumntype{Y}{>{\raggedright\arraybackslash}X}
\newcolumntype{P}[1]{>{\raggedright\arraybackslash}p{#1}}

\title{\field{Project Title: Identify the Problem and Core Idea}}
\author{\IEEEauthorblockN{\field{Student 1 (UNI)}, \field{Student 2 (UNI)}, \field{Student 3 (UNI)}}
\IEEEauthorblockA{COMS 6998: AI-Native Computing $\mid$ Fall 2026\\
Project Proposal $\mid$ Due October 5, 2026, 11:59 PM ET}}

\begin{document}
\maketitle
\thispagestyle{plain}
\pagestyle{plain}

\noindent\textbf{Track / topic:} \field{A: Computing for AI, or B: AI for Computing; project topic.}

\smallskip
\noindent\textbf{Format.} Write \textbf{two pages}, excluding references.
Use the five sections below
and replace the guidance and bracketed fields with your project content.

\section{Problem and Motivation}
\textbf{Problem.} Describe the problem you want to address in a specific
workload, application, or computing system.

\textbf{Motivation.} Explain why it's important and how addressing it
could benefit the target application or system.

\textbf{Research question.} State the main question your project will
explore and what you aim to achieve.

\section{Related Work}
Briefly discuss a few relevant papers or systems and explain how your
project builds on or differs from them. You can highlight an opportunity
for improvement or a question worth exploring. Use citations such
as~\cite{dao2022flashattention}; replace this example with your own references.

\section{Proposed Approach}
Describe your main idea and how you plan to implement it. You can outline
the key components, the tools or codebases you will use, and the expected
contribution. A diagram can help explain the approach.

\section{Evaluation Plan}
Outline how you plan to evaluate your idea. You can mention potential
workloads or datasets, methods to compare against, and useful metrics
such as task quality, latency, or cost. An initial plan is sufficient;
experiment details can be refined as the project develops.

\section{Timeline and Team Responsibilities}
Use the course dates below to outline your planned work and expected
deliverables for each stage.

\begin{center}
\begin{tabularx}{\columnwidth}{@{}P{0.23\columnwidth}Y@{}}
\toprule
\textbf{Date} & \textbf{Planned work and deliverables} \\
\midrule
Oct 19 & \field{Your plan.} \\
Nov 2 & \field{Your plan.} \\
Nov 23 & \field{Your plan.} \\
Dec 7 & \field{Your plan.} \\
Dec 18 & \field{Your plan.} \\
\bottomrule
\end{tabularx}
\end{center}

\textbf{Presentations.} Midterm presentation: Nov~6. Final poster session: Dec~11.

\textbf{Team roles.} Briefly describe the overall division of work within
your team.

\bibliographystyle{IEEEtran}
\bibliography{references}

\end{document}
